% !PRINTED 285
\weberpage{285}

If we now, as is customary in number theory (\gesperrt{Gauss}, Disquisitiones arithmeticae, art. 31), understand by the symbol
\[
\frac{a}{b} \ (\mathrm{mod}\ p)
\]
an \gesperrt{integer} [or class of numbers (mod $p$)] which, multiplied by $b$, leaves the residue $a$ on division by $p$, then, provided that $\mu'$ and $\nu'$ are not divisible by $p$, we can also write the congruence (2) that defines a row as follows:
\weq{3}{\frac{\mu}{\mu'} \equiv \frac{\nu}{\nu'} \ (\mathrm{mod}\ p),}
and we are thus naturally led to denote this row $R$ by the value of the ratio
\weq{4}{\frac{\mu}{\mu'} \equiv z \ (\mathrm{mod}\ p),}
which can be congruent to any one of the numbers $0, 1, \ldots, p-1$ and which is invariable for a whole row. At first one row here remains undesignated, namely that in which $\mu'$, and consequently all $\nu'$, are divisible by $p$; but this row too is subsumed very well under the general designation if, in the case $\mu' \equiv 0 \ (\mathrm{mod}\ p)$, we put
\weq{5}{\frac{\mu}{\mu'} \equiv \infty \ (\mathrm{mod}\ p)}
so that we thus obtain a further $(p+1)$th row $R_\infty$. The totality of the rows, whose number amounts to $p+1$, is accordingly to be denoted by
\weq{6}{R_\infty, R_0, R_1, \ldots, R_{p-1}.}

Correspondingly, for the associated roots of a transformation equation the designations
\weq{7}{v_\infty, v_0, v_1, \ldots, v_{p-1}}
will have to be used.

If, for example, we take the invariant equation (§ 69) as our basis, then [according to the determinations of § 68, (10) concerning the numbers $a$, $c$, $\partial$] we must put
\weq{8}{v_\infty = j(p\,\omega), \quad v_z = j\left(\frac{z+\omega}{p}\right)}
and likewise, if any other transformation equation is chosen.
